% Archivo compartido: no compilar directamente.
\ifdefined\documentoSemanal
\else
\usepackage{tikz}
\usepackage{estilo_apuntes}
\usetikzlibrary{babel}
\encabezado{Semana 2 · Clase 2}{Material de clase}
\begin{document}
\fi

\ifdefined\documentoSemanal
\else
\tituloclase{Noción intuitiva de límite}{Tablas, gráficas y sucesiones}{Semana 2 · Clase 2 · 2 horas}

\begin{objetivos}
\begin{itemize}
  \item Interpretar $\lim_{x\to a}f(x)$ mediante valores próximos a $a$.
  \item Contrastar evidencia tabular, gráfica y obtenida con sucesiones.
  \item Distinguir límite, límites laterales y valor de la función.
  \item Comprender la unicidad y el principio de funciones localmente iguales.
\end{itemize}
\end{objetivos}

\begin{resumen}
El límite describe el comportamiento de $f(x)$ cuando $x$ se aproxima a $a$, sin exigir que $x=a$. Por ello puede existir aunque $f(a)$ no exista o sea distinto del límite. Las tablas y gráficas sugieren el valor; las sucesiones precisan qué significa aproximarse por cualquier camino.
\end{resumen}
\fi

\section*{Aproximación a un punto}

\begin{definicion}[noción intuitiva de límite]
Escribimos
\[
\lim_{x\to a}f(x)=L
\]
cuando los valores $f(x)$ pueden hacerse tan próximos a $L$ como se desee al tomar $x$ suficientemente próximo a $a$, con $x\neq a$. El límite estudia valores en un entorno perforado de $a$.
\end{definicion}

\begin{nota}
La expresión $x\to a$ no significa que $x$ llegue a ser igual a $a$. Las afirmaciones $\lim_{x\to a}f(x)=L$ y $f(a)=L$ son diferentes y ninguna implica por sí sola la otra.
\end{nota}

\begin{ejemplo}[exploración mediante una tabla]
Para $x\neq1$, considere
\[
f(x)=\frac{x^2-1}{x-1}.
\]
Complete o verifique la tabla y conjeture el límite cuando $x\to1$.
\[
\begin{array}{c|cccccc}
x&0.9&0.99&0.999&1.001&1.01&1.1\\ \hline
f(x)&1.9&1.99&1.999&2.001&2.01&2.1
\end{array}
\]
\end{ejemplo}
\ifmostrarSoluciones
\begin{solucion}
Los valores se aproximan a $2$ desde ambos lados. Como $x^2-1=(x-1)(x+1)$,
\[
f(x)=x+1\quad(x\neq1),
\]
y la evidencia sugiere $\lim_{x\to1}f(x)=2$, aunque $f(1)$ no está definida.
\end{solucion}
\else\vspace{12mm}\fi

\begin{center}
\begin{tikzpicture}[x=1.05cm,y=0.8cm,every node/.style={font=\scriptsize}]
  \draw[Gris!70] (-0.4,0)--(3.1,0) node[right] {$x$};
  \draw[Gris!70] (0,-0.3)--(0,3.4) node[above] {$y$};
  \draw[Azul,very thick,domain=-0.1:2.8,samples=2] plot (\x,{\x+1});
  \draw[fill=white,draw=Rojo,thick] (1,2) circle (2.2pt);
  \fill[Rojo] (1,3) circle (1.8pt);
  \draw[Gris,dashed] (1,0)--(1,3);
  \node[below] at (1,0) {$a$};
  \node[left] at (0,2) {$L$};
  \node[right] at (1,2) {valor al que se aproxima};
  \node[right] at (1,3) {$f(a)$};
\end{tikzpicture}
\end{center}

\begin{definicion}[límites laterales]
El \textbf{límite por la izquierda}, $\lim_{x\to a^-}f(x)$, utiliza valores $x<a$. El \textbf{límite por la derecha}, $\lim_{x\to a^+}f(x)$, utiliza valores $x>a$.
\end{definicion}

\begin{teorema}[criterio de los límites laterales]
El límite bilateral existe y vale $L$ si y solo si existen los dos límites laterales y coinciden:
\[
\lim_{x\to a}f(x)=L
\iff
\lim_{x\to a^-}f(x)=L=\lim_{x\to a^+}f(x).
\]
Si los límites laterales son distintos, el límite bilateral no existe.
\end{teorema}

\begin{ejemplo}[salto]
Sea $s(x)=-1$ para $x<0$ y $s(x)=2$ para $x\geq0$. Determine sus límites laterales y decida si existe $\lim_{x\to0}s(x)$.
\end{ejemplo}
\ifmostrarSoluciones
\begin{solucion}
\[
\lim_{x\to0^-}s(x)=-1,\qquad
\lim_{x\to0^+}s(x)=2.
\]
Como son distintos, $\lim_{x\to0}s(x)$ no existe, aunque $s(0)=2$ sí está definido.
\end{solucion}
\else\vspace{13mm}\fi

\section*{Interpretación mediante sucesiones}

\begin{proposicion}[criterio secuencial intuitivo]
La afirmación $\lim_{x\to a}f(x)=L$ exige que, para toda sucesión $(x_n)$ del dominio tal que
\[
x_n\neq a\quad\text{y}\quad x_n\to a,
\]
se tenga $f(x_n)\to L$. Para refutar la existencia del límite basta encontrar dos sucesiones que se aproximen a $a$ y cuyas imágenes tengan límites diferentes.
\end{proposicion}

\begin{ejemplo}[dos caminos hacia el mismo punto]
Use $x_n=1/n$ y $y_n=-1/n$ para estudiar el límite en $0$ de la función signo, definida como $1$ para entradas positivas y $-1$ para entradas negativas.
\end{ejemplo}
\ifmostrarSoluciones
\begin{solucion}
Ambas sucesiones tienden a $0$, pero $\operatorname{sgn}(x_n)=1$ y $\operatorname{sgn}(y_n)=-1$ para todo $n$. Las imágenes tienen límites distintos; por tanto, el límite en $0$ no existe.
\end{solucion}
\else\vspace{13mm}\fi

\section*{Unicidad y comportamiento local}

\begin{teorema}[unicidad del límite]
Si $\lim_{x\to a}f(x)$ existe, su valor es único. En consecuencia, una misma función no puede aproximarse simultáneamente a dos números distintos cuando $x\to a$.
\end{teorema}

\begin{definicion}[funciones localmente iguales]
Las funciones $f$ y $g$ son \textbf{localmente iguales alrededor de $a$} si existe un intervalo abierto $I$ que contiene a $a$ tal que
\[
f(x)=g(x)\qquad\text{para todo }x\in I\setminus\{a\}.
\]
Pueden diferir en $a$ o fuera de ese intervalo.
\end{definicion}

\begin{proposicion}[principio de sustitución local]
Si $f$ y $g$ son localmente iguales alrededor de $a$, entonces
\[
\lim_{x\to a}f(x)=L\iff\lim_{x\to a}g(x)=L.
\]
El límite depende del comportamiento próximo al punto, no del valor asignado exactamente en él ni de lo que ocurra lejos de él.
\end{proposicion}

\begin{ejemplo}[cambio en un solo punto]
Considere $g(x)=x+1$ y defina $f(x)=x+1$ cuando $x\neq1$, pero $f(1)=7$. Compare $f(1)$, $g(1)$ y los límites de ambas funciones cuando $x\to1$.
\end{ejemplo}
\ifmostrarSoluciones
\begin{solucion}
$f(1)=7$ y $g(1)=2$, pero las funciones coinciden para todo $x\neq1$. Por el principio de sustitución local,
\[
\lim_{x\to1}f(x)=\lim_{x\to1}g(x)=2.
\]
\end{solucion}
\else\vspace{12mm}\fi

\begin{ejercicios}[noción intuitiva de límite]
\begin{enumerate}
  \item Construya una tabla para estimar $\lim_{x\to2}(x^2+1)$.
  \item Estime $\lim_{x\to3}\dfrac{x^2-9}{x-3}$ con tres valores a cada lado de $3$.
  \item Dibuje una función con $f(2)=4$ y $\lim_{x\to2}f(x)=1$.
  \item Decida si existe el límite en $0$ de $f(x)=|x|/x$.
  \item Use dos sucesiones para justificar su respuesta en el ejercicio anterior.
  \item Proponga dos funciones localmente iguales alrededor de $2$ que tengan valores diferentes en $2$.
  \item ¿Puede una función tener límite $3$ y límite $5$ en el mismo punto? Justifique.
\end{enumerate}
\end{ejercicios}

\ifmostrarSoluciones
\begin{solucion}[de los ejercicios]
\begin{enumerate}
  \item Los valores se aproximan a $5$.
  \item Para $x\neq3$, el cociente es $x+3$; la tabla se aproxima a $6$.
  \item Por ejemplo, $f(x)=1$ si $x\neq2$ y $f(2)=4$.
  \item No existe: los límites laterales son $-1$ y $1$.
  \item Con $x_n=1/n$ y $y_n=-1/n$, las imágenes son $1$ y $-1$, respectivamente.
  \item Por ejemplo, $f(x)=x^2$ y $g(x)=x^2$ para $x\neq2$, con $g(2)=0$.
  \item No, por la unicidad del límite.
\end{enumerate}
\end{solucion}
\fi

\ifdefined\documentoSemanal
\else
\fuentes{\textbf{Para ampliar:} \emph{Cálculo en una variable: notas de clase 2018-B}, apartado de límites; J. Stewart, \emph{Cálculo de una variable: trascendentes tempranas}, §2.2; y G. Strang--E. Herman, \emph{Cálculo, volumen 1}, OpenStax, §2.2.}
\end{document}
\fi
